\DocumentMetadata{lang=es-CL,pdfversion=2.0,pdfstandard=ua-2,tagging=on,tagging-setup={math/setup={mathml-SE, mathml-AF}},}
\documentclass{article}
\usepackage[spanish]{babel}
\ifdefined\Uchar
  \usepackage{unicode-math}
\fi

\title{babel-spanish tagging test}

\begin{document}

"u "U "c "C "rr "RR "y "/ \'{\i}

$\acute{\dotlessi}$

\section{Comillas Latinas o Angulares}
Las comillas latinas (« y ») son las preferidas por la RAE. Se activan con los
atajos \verb+"<+ y \verb+">+:

\begin{quotation}
El autor afirmó: "<Esto es una cita textual con shorthands">.
\end{quotation}

\section{Gestión de Guiones e Hifenación}
Babel introduce variantes del carácter guión para controlar cómo el motor de
\LaTeX{} corta las palabras al final de las líneas.

\begin{itemize}
    \item \textbf{Guión compuesto libre:}  (\verb+"-+) Permite que el resto
    de la palabra se divida si es necesario. \par
    Ejemplo (físico"-químico): físico"-químico

    \item \textbf{Guión compuesto rígido:} (\verb+"=+) Introduce un guión
    pero prohíbe taxativamente separar las palabras contiguas. \par
    Ejemplo (chileno"=argentina): chileno"=argentina

    \item \textbf{Punto de rotura invisible:} (\verb+""+) Permite romper la
    línea en ese punto exacto sin añadir ningún guión visual. \par
    Ejemplo (kilómetros""/""hora): kilómetros""/""hora
\end{itemize}

\section{Signos de Interrogación y Exclamación}
Los atajos introducen un microespacio tipográfico elástico automático:

\begin{itemize}
    \item \textbf{Interrogación:} (\verb+"?+) "?Cómo funciona esto?
    \item \textbf{Exclamación:} (\verb+"!+) "!Qué excelente resultado!
\end{itemize}

\section{Comillas Curvas Auxiliares}
Para anidar comillas dentro de las latinas, se usan los atajos \verb+"`+
y \verb+"'+:

\begin{quotation}
<<El texto decía "`palabra"' dentro de comillas>>.
\end{quotation}

\section{Abreviaturas Estándar con Letras Voladas}
La RAE exige que las letras voladas vayan precedidas de un punto. Para generar
abreviaturas correctas en el ecosistema, se utiliza la macro nativa \verb+\sptext+:

\begin{itemize}
    \item \textbf{Abreviaturas comunes:} n\sptext{o}~1, art\sptext{o}~4, pág.~42.
    \item \textbf{Ordinales y títulos:} 1\sptext{er} (Primer), 1\sptext{o} (Primero), Prof\sptext{a}.
\end{itemize}

\section{Operadores Matemáticos en Español}
Babel redefine automáticamente los comandos de funciones trigonométricas y
logarítmicas para que se impriman en castellano dentro del modo matemático:

\begin{itemize}
    \item \textbf{Seno:} $\sen(x)$ (en lugar del inglés \emph{sin}).
    \item \textbf{Tangente:} $\tan(x)$ o también usando $\tg(x)$.
    \item \textbf{Arcotangente:} $\arctg(x)$.
    \item \textbf{Logaritmo neperiano:} $\ln(x)$.
\end{itemize}

\end{document}